% ============================================
% DOC-12
% Gerado de molambudos.md — build_miolo.py
% ============================================
\fragtitle{doc}{D}{DOC}{12}{Diário de Oliveira II}
\markboth{DOC-12 --- Diário de Oliveira II}{}

\noindent \textit{A letra muda}

\noindent \textbf{Dr. Heitor Oliveira --- Caderno Pessoal}
\textbf{20 de abril de 1979}

\vspace{0.5em}\noindent\rule{\textwidth}{0.5pt}\vspace{0.5em}

\noindent Dia 7.

\noindent Minha letra está mudando.

\noindent Percebi hoje de manhã, quando fiz anotações sobre um paciente. Olhei o prontuário e não reconheci a caligrafia. Parecia mais... torta. Infantil. Como a letra de alguém que aprendeu a escrever sozinho.

\noindent Como a letra de Joaquim.

\noindent Peguei o diário dele --- sim, ainda está comigo, não consigo me livrar dele --- e comparei. A similaridade é impressionante. O \textquotedblleft{}J\textquotedblright{} maiúsculo, o \textquotedblleft{}a\textquotedblright{} minúsculo, o traço final das palavras.

\noindent Estou copiando a letra dele sem perceber.

\noindent Ou ele está copiando a minha.

\noindent Ou a mesma mão está escrevendo as duas coisas.

\vspace{0.5em}\noindent\rule{\textwidth}{0.5pt}\vspace{0.5em}

\noindent Dia 8.

\noindent Hoje tentei me livrar do diário.

\noindent Coloquei dentro de uma sacola plástica. Amarrei. Saí de carro até um ponto de ônibus na BR-040, a 15 quilômetros de Barbacena. Joguei a sacola numa lixeira pública.

\noindent Voltei para casa.

\noindent O diário estava na mesa da cozinha.

\noindent Sem sacola. Sem amarras. Apenas o diário, aberto na página onde eu tinha parado de ler.

\noindent Minha esposa disse: \textquotedblleft{}Alguém trouxe?\textquotedblright{}

\noindent Eu disse: \textquotedblleft{}Não.\textquotedblright{}

\noindent Ela disse: \textquotedblleft{}Então como ele voltou?\textquotedblright{}

\noindent Eu não respondi.

\noindent Peguei o diário. Guardei no bolso do casaco. Dessa vez, não vou tentar me livrar. Não adianta.

\noindent Ele quer ser lido.
Ele \textit{precisa} ser lido.
E eu preciso ler.

\vspace{0.5em}\noindent\rule{\textwidth}{0.5pt}\vspace{0.5em}

\noindent Dia 9.

\noindent Li mais 50 páginas.

\noindent Joaquim descreve o encontro na vala em detalhes que nenhum delírio poderia produzir. A criatura falou com ele. Fez um trato. Deu a fome em troca de companhia.

\noindent Há uma passagem que me gelou:

\noindent \textit{\textquotedblleft{}Ela disse: 'Um dia, um homem de jaleco branco vai ler este diário. Ele vai negar. Ele vai ter medo. Ele vai tentar fugir. Mas ele vai voltar. Porque a curiosidade é maior que o medo. E quando ele voltar, eu vou estar esperando.'\textquotedblright{}}

\noindent \textit{\textquotedblleft{}Perguntei: 'Como você sabe?'\textquotedblright{}}

\noindent \textit{\textquotedblleft{}Ela disse: 'Porque ele já esteve aqui antes. Em outra vida. Em outro corpo. Ele sempre volta. Desde o início.'\textquotedblright{}}

\noindent Desde o início.

\noindent Quantos leitores vieram antes de mim?

\noindent Quantos médicos, pesquisadores, curiosos abriram este diário e sentiram o cheiro, e viram o olho, e começaram a sonhar?

\noindent Onde estão eles agora?

\vspace{0.5em}\noindent\rule{\textwidth}{0.5pt}\vspace{0.5em}

\noindent Dia 10.

\noindent Pesquisei.

\noindent O diário foi mencionado em três artigos acadêmicos nos anos seguintes. Todos os autores tiveram \textquotedblleft{}problemas de saúde\textquotedblright{} antes de publicar. Um morreu de câncer. Outro teve um colapso nervoso. O terceiro --- uma psiquiatra chamada Dra. Verônica Tavares --- desapareceu em 1985.

\noindent O diário foi encontrado na casa dela. Junto com um bilhete:

\noindent \textit{\textquotedblleft{}O ciclo não pode ser quebrado. Só passado adiante. Devolvo ao arquivo para que outro leia. Desculpe.\textquotedblright{}}

\noindent Ela passou adiante.

\noindent E eu peguei.

\vspace{0.5em}\noindent\rule{\textwidth}{0.5pt}\vspace{0.5em}

\noindent Dia 11.

\noindent 3:14 da manhã.

\noindent Acordei com o olho arregalado.

\noindent Não é metáfora. O olho direito estava aberto. Não conseguia fechar. A pálpebra estava dura.

\noindent Fui ao espelho.

\noindent O olho direito está amarelo.

\noindent Não vou ao médico. Não vou contar a ninguém. Já li o suficiente para saber o que está acontecendo.

\noindent A criatura está passando para mim.

\noindent Joaquim morreu.
Eu li o diário.
Agora ela está em mim.

\noindent Meu olho está amarelo.
Meu cheiro está doce.
Minha letra está igual à dele.

\noindent Estou me tornando o paciente 1.261.

\vspace{0.5em}\noindent\rule{\textwidth}{0.5pt}\vspace{0.5em}

\noindent Dia 12.

\noindent Decidi.

\noindent Vou escrever meu próprio diário. Vou registrar o que está acontecendo comigo. Para que, quando alguém ler, saiba.

\noindent O ciclo é uma corrente de leitores.

\noindent Joaquim → Oliveira → ? → ? → ?

\noindent Alguém vai ler isto depois de mim.

\noindent Se você está lendo: sinto muito.

\noindent Não há como parar.

\noindent O diagnóstico não é metáfora.

\noindent A criatura é real.

\noindent Eu vi o olho.

\noindent Eu senti o cheiro.

\noindent Eu sou o próximo.

\noindent \textit{--- Heitor Oliveira, 29 de abril de 1979}
\textit{--- Paciente 1.261}

\vspace{0.5em}
\noindent\textit{\textbf{↪ Rotas:} \rota{DOC-13} \rotasep \rota{CONT-01}}